\documentclass[11pt]{article}
\usepackage[margin=1in]{geometry}
\usepackage{amsmath,amssymb,booktabs,hyperref,natbib}
\usepackage[T1]{fontenc}

% =================================================================
% GAPS.tex -- Validation notes on the Problem Formulation.
%
% Supersedes GAPS.md (kept as markdown was the wrong format for a document
% whose whole content is thesis equations, cross-references, and a table
% meant to be pasted into the thesis's own appendix -- it belongs typeset
% alongside the thesis it discusses, not as a loose Markdown file).
%
% This file stands alone (compiles with `pdflatex GAPS.tex`) but its
% equation numbers/labels intentionally mirror problem_formulation_fixed.tex,
% so \eqref calls below refer to that file's labels, not local equations,
% wherever they cite eq:xxx from Section IV. Re-numbering is not attempted
% here; this is a companion document, not a subfile of the thesis itself.
%
% Two parts:
%   Part I  -- internal-consistency review of problem_formulation.tex
%              ([G1]-[G13], matching the inline tags in the fixed .tex)
%   Part II -- comparison against Knippenberg, Holenderski & Menkovski (2021),
%              the closest prior-art paper this section's environment and
%              observation model is adapted from ([K1]-[K6])
%   Appendix -- the notation-disambiguation table, appendix-ready LaTeX to
%              paste into the thesis's own appendix chapter (not present in
%              this repo) rather than left inline in Section IV, per the
%              Problem Formulation / Method / Implementation contract in
%              Method_and_Implementation.md.
% =================================================================

\title{Validation Notes on the Problem Formulation}
\author{}
\date{}

\begin{document}
\maketitle

\tableofcontents

\section*{How to use this document}
\texttt{problem\_formulation\_fixed.tex} is a complete, label-preserving
drop-in revision of \texttt{problem\_formulation.tex}. Every fix there is
tagged inline (\texttt{\% [Gxx]} or \texttt{\% [Kxx]}); this document is the
rationale for each tag, split into two parts by where the finding came from:
Part~\ref{sec:internal} was found by checking the section against itself
(does every symbol get defined, do the claimed properties actually follow);
Part~\ref{sec:knippenberg} was found by checking it against the paper its
environment and observation model is adapted from. Severity key used
throughout: \textbf{Error} = mathematically wrong as stated; \textbf{Gap} =
not wrong, but undefined/underspecified where the section otherwise commits
to full formality; \textbf{Deviation} = a considered, defensible difference
from a source, not necessarily a defect; \textbf{Clarity} = correct and
complete, but hard to read.

\part*{}
\section{Internal-consistency review}
\label{sec:internal}

\subsection*{[Error] The potential-function claim (target-distance annotation)}
\textbf{Original text:} ``The annotation is a potential: its difference along
an edge, $\eta_i(w,t) - \eta_i(v,t)$, is negative exactly on shortest-path
moves.''

\textbf{Why it's wrong:} by the triangle inequality on $d_G$,
\[
  \eta_i(w,t) - \eta_i(v,t) \;\ge\; -c(v,w) \qquad \text{for every } (v,w)\in E,
\]
with \emph{equality iff} $(v,w)$ lies on a shortest path from $v$ to
$q_i(t)$. ``Negative'' is necessary for a shortest-path move but not
sufficient: a move can strictly satisfy the inequality (decrease $\eta_i$ by
less than $c(v,w)$) while not lying on any shortest path. Concretely:
$c(v,w)=1$, $d_G(v,q)=5$, and some other route makes $d_G(w,q)=4.5$ -- the
difference is $-0.5$ (negative), but $(v,w)$ is not on a shortest path (that
would require $d_G(w,q)=4$ exactly).

\textbf{Fix:} restated as the inequality with the correct equality condition
(\texttt{eq:potentialbound} in the fixed .tex); surrounding prose now says a
negative difference indicates \emph{progress}, not \emph{optimality}. This is
the one place a formally stated property was actually false, not merely
underspecified.

\subsection*{[G1] \texorpdfstring{$V_{mov}$}{Vmov} vs. \texorpdfstring{$V_{str}$}{Vstr}/\texorpdfstring{$V_{del}$}{Vdel} -- can an agent ever reach a pickup or delivery vertex?}
Agents only ever occupy $\ell_i(t) \in V_{mov}$. The original text declares
$V_{str}, V_{del} \subseteq V$ (not $\subseteq V_{mov}$) while calling all
four roles a ``partition'' -- which $V_{ep}\subseteq V_{mov}$ already
contradicts, since a partition requires disjoint blocks. If $V_{str}$ and
$V_{del}$ are disjoint from $V_{mov}$, as ``partition'' suggests, an agent can
never occupy a pickup or delivery vertex, leaving pickup/delivery physically
undefined. \textbf{Fix:} $V_{str}, V_{del} \subseteq V_{mov}$ now stated
explicitly (they must be occupiable), ``partition'' replaced by ``label,''
and pairwise disjointness explicitly not assumed.

\subsection*{[G2] The task tuple is missing the SKU it names}
$\tau_j = (r_j, s_j, g_j)$ in the original, but the coupling condition
immediately after it, $x_{r_j}(k_j, s_j) > 0$, uses $k_j$ -- the requested
SKU -- which is never introduced anywhere. The thesis's central equation
rests on an undefined symbol. \textbf{Fix:} \texttt{eq:task} is now a
quadruple $\tau_j = (r_j, s_j, g_j, k_j)$, $k_j \in \mathcal{K}$.

\subsection*{[G3] \texorpdfstring{$\Pi$}{Pi} and \texorpdfstring{$\mathcal{F}$}{F} are never defined as sets}
The formal decision problem optimises over $F \in \mathcal{F}$, $\pi \in
\Pi$, but neither space's domain/codomain is ever written down -- only
\emph{what information} each may condition on (Table~I) is described.
\textbf{Fix:} $\mathcal{F}$ defined as the set of measurable functions
$F : \mathcal{X} \times \mathbb{R}_{\ge0}^{|\mathcal{K}|} \times
\mathbb{R}_{\ge0}^{|E|} \times \mathbb{R}_{\ge0}^{|E|} \to \mathcal{X}$;
$\Pi$ defined as the set of pairs $\pi=(\pi^{\text{assign}},
\pi^{\text{route}})$ of an assignment map and a routing map with stated
domains, constrained only by Table~I. A related loose end fixed in the same
place: eq.~8 ($x_t \in \arg\min$) can be set-valued while eq.~7 ($x_t =
F(\ldots)$) claims a function; the fix adds a tie-breaking-rule requirement
so eq.~7 is well-defined.

\subsection*{[G4] No initial storage configuration}
$x_t = F(x_{t-\Delta}, \ldots)$ needs $x_0$ to bootstrap it at $t=\Delta$;
$x_0$ is never introduced. \textbf{Fix:} $x_0 \in \mathcal{X}$ stated as
supplied with the problem instance.

\subsection*{[G5] \texorpdfstring{$\mathcal{Q}_t, \mathcal{B}_t, \mathcal{C}_t$}{Qt, Bt, Ct} are only described narratively}
Disjointness is asserted (``a property asserted at runtime'') rather than
derived, unlike every other object in the section. \textbf{Fix:} introduces
the assignment time $y_j$ ($r_j \le y_j \le d_j$) and defines
\[
\mathcal{Q}_t = \{\tau_j : r_j \le t < y_j\}, \quad
\mathcal{B}_t = \{\tau_j : y_j \le t < d_j\}, \quad
\mathcal{C}_t = \{\tau_j : d_j \le t\}.
\]
Disjointness now follows from the three intervals partitioning
$[r_j,\infty)$, rather than being separately asserted.

\subsection*{[G6] The congestion feature \texorpdfstring{$\delta_i(t)$}{delta\_i(t)} is defined but never wired in}
$\delta_i(t)$ appears in the Learning Problem subsection, but $o_i(t)$ was
fixed earlier and never amended, and $\delta_i(t)$ is never connected to
$R_i$ either. \textbf{Fix:} moved \texttt{eq:congestion} into the
Observations subsection, right after $\eta_i$, and added $\delta_i(t)$ as a
third component of \texttt{eq:observation}.

\subsection*{[G7] \texorpdfstring{$W_T$}{WT} has no formula and an ambiguous relationship to \texorpdfstring{$\hat{w}_t(e)$}{what(e)}}
$W_T$ appears only inside the system-level criterion, described as having an
operational definition deferred to Implementation, with zero conceptual
description here -- and a name close enough to the already-used per-edge
estimate $\hat{w}_t(e)$ that the two could be mistaken for the same object.

\textbf{First-pass fix (superseded below):} an initial revision gave
$W_T = \frac{1}{|\mathcal{C}_T|}\sum_{\tau_j \in \mathcal{C}_T}(y_j - r_j)$,
mean pre-assignment queueing delay, using $y_j$ from [G5]. This was internally
consistent but, once \texttt{Method\_and\_Implementation.md}'s operational
suggestion for $W_T$ (``timesteps in which an agent waited or was overridden
while holding an active task'') was checked against it, turned out to describe
a \emph{different} quantity -- and queueing delay is largely redundant with the
backlog term $B_T = |\mathcal{Q}_T|+|\mathcal{B}_T|$ already in
\texttt{eq:functional}, the same double-counting concern the section already
raises for throughput vs.\ $\bar\zeta_T$.

\textbf{Current fix:} $W_T$ now measures time lost \emph{after} assignment
instead, which is not redundant with $B_T$: for agent $i$ serving $\tau_j$ over
$t\in[y_j,d_j)$, let $\omega_j(t)\in\{0,1\}$ indicate agent $i$ is blocked at
$t$ (chose wait, or had a move overridden into a wait by the conflict-resolution
operator, Section~\ref{sec:method-mapping}/\texttt{sec:method:resolution}).
Then
\[
W_T = \frac{1}{|\mathcal{C}_T|} \sum_{\tau_j\in\mathcal{C}_T}
      \sum_{t=y_j}^{d_j-1} \omega_j(t).
\]
This version also gives the resolution operator's ``overridden'' flag (needed
anyway for the feasibility mechanism) a direct second use as $W_T$'s data
source, rather than requiring separate instrumentation. $W_T$ remains distinct
from $\hat w_t(e)$ for the same reason as before: one is a realised scoring
statistic, the other a forward-looking estimate feeding $F$.

\subsection*{[G8] Entropy is undefined at the boundary}
\texttt{eq:entropy} is stated ``for $|E_T^+| > 1$'' only; at
$|E_T^+|\in\{0,1\}$ the denominator $\log|E_T^+|=0$. \textbf{Fix:}
$|E_T^+|=1$ defined as maximally concentrated by convention ($H_T := 0$,
$C_T := 1$); $|E_T^+|=0$ reported separately, like $|\mathcal{C}_T|=0$.

\subsection*{[G9] Are constraints (4)--(5) actually constraining anything in the decision problem?}
Section~B already states vertex/swap-conflict feasibility is enforced by a
mechanism applied regardless of which $\pi$ is chosen. If so, every $\pi$
satisfies them by construction, and listing them in the $\arg\min$ over
$(F,\pi)$ doesn't prune the search space. \textbf{Fix:} one clarifying
paragraph: these constraints are retained to state explicitly that a
solution is a claim about \emph{feasible realised trajectories}, unlike
\texttt{eq:feasiblestorage} and \texttt{eq:coupling}, which genuinely do
restrict $F$ and the realised task stream.

\subsection*{[G10] Is \texorpdfstring{$G$}{G} directed, and is it connected?}
One-way aisles imply $G$ must be directed, but this is never stated; $d_G$ is
used pervasively as if always finite, with no connectivity assumption
stated. \textbf{Fix:} $G$ explicitly stated as directed; a finiteness
assumption added (tightened further in Part~\ref{sec:knippenberg}, [K1]).

\subsection*{[Gap] \texorpdfstring{$\mathcal{U}_i$}{Ui} (POSG action space) vs. \texorpdfstring{$\mathcal{U}(v)$}{U(v)} (vertex-local action set)}
\texttt{eq:actions} is explicitly vertex-dependent, but the POSG tuple lists
a fixed per-agent $\mathcal{U}_i$, with no statement of how one becomes the
other. \textbf{Fix:} $\mathcal{U}_i := \bigcup_{v\in V_{mov}}\mathcal{U}(v)$,
with only $\mathcal{U}(\ell_i(t))$ enabled at a given state (the rest masked
to probability zero).

\subsection*{[Gap] Does \texorpdfstring{$J(\theta)$}{J(theta)} apply to all three controller architectures?}
$\theta$ is only introduced for the decentralised architecture, but
\texttt{eq:objective} is presented as though general. \textbf{Fix:} one
sentence: it applies to whichever regime is realised as a learned
controller, with $\theta$ ranging over that regime's own parameters.

\subsection*{[G13] Is the conflict-resolution mechanism itself deterministic?}
Section~J claims execution is deterministic with $P$ stochastic only through
order arrivals, but never confirms the yield-resolution mechanism is itself
deterministic -- if it used randomised tie-breaking, that would be a second
stochasticity source. \textbf{Fix:} one sentence confirming it is
deterministic given the state.

\subsection*{[G11 / Clarity] Symbol collisions: five unrelated concepts spelled with the letter ``C,'' three with ``Q,'' two with ``O''}
Not an error, but slows careful reading (eq.~21 has both $\bar c_T$ and
$\mathcal{C}_T$ in the same line). \textbf{Fix:} rather than renaming
established symbols this late -- Method/Implementation may already reference
them, and a rename should touch the whole thesis at once, not one section in
isolation -- a disambiguation table is provided in
Appendix~\ref{app:notation} of \emph{this} document, ready to paste into the
thesis's own appendix chapter, with only a one-line pointer left inline in
Section~IV (per \texttt{Method\_and\_Implementation.md}'s own rule that a
full notation table is appendix material).

\subsection*{Left open / not changed}
\begin{itemize}
  \item \textbf{$m$ used before being defined} (Section~A uses it,
    Section~B defines $m=|\mathcal{A}|$) -- fixed with a one-line forward
    pointer rather than reordering subsections.
  \item \textbf{A full rename of the C/Q/O-family symbols} -- flagged via
    the appendix table rather than executed; out of scope for a
    Problem-Formulation-only fix.
  \item \textbf{Unit/scale normalisation across the five terms of
    $\mathcal{J}(F,\pi)$} -- summing time, cost, task-count and
    dimensionless quantities under weights $q_\bullet$ with no stated
    calibration scheme is a Method-level concern (how $q_\bullet$ are
    chosen), not a Problem-Formulation error, but worth having an answer
    ready for: linear scalarisation across unlike units needs either
    normalisation or explicit justification for the weights.
  \item \textbf{Section~H's $u_t\sim\pi(\cdot\mid s_t)$ vs.\ the
    decentralised $\pi_\theta(\cdot\mid o_i(t))$} -- already reconciled
    informally (``with the corresponding product of local policies under
    decentralised execution''); not tightened further.
\end{itemize}

\section{Comparison against Knippenberg, Holenderski \& Menkovski (2021)}
\label{sec:knippenberg}

\citet{knippenberg2021} (``Time-Constrained Multi-Agent Path Finding in
Non-Lattice Graphs with Deep Reinforcement Learning,'' ACML 2021) is the
closest prior-art paper to this thesis's environment and observation model:
a directed, non-lattice graph; a local observation of fixed depth $d$
centred on the agent; a per-node distance-to-goal annotation; a shared
policy over a wait-or-move action space. The thesis's problem formulation is
best read as taking that architecture and (a) embedding it inside a
lifelong/online MAPD setting with continuous task arrivals rather than
\citeauthor{knippenberg2021}'s batch/episodic one-job-per-agent setting, and
(b) coupling it to a new second layer -- storage placement -- that has no
counterpart in \citeauthor{knippenberg2021} at all. What follows is organised
by what carried over unchanged, what was correctly adapted, what is a new
extension, and one place the comparison surfaces an actual gap.

\subsection{Carried over directly (and appropriately)}
\begin{itemize}
  \item \textbf{Directed, non-lattice graph.} Identical framing in both.
  \item \textbf{Local observation of fixed depth $d$.}
    \citeauthor{knippenberg2021}'s observation is ``centered on the agent's
    current location'' with size ``determined by a global depth parameter
    $d$''; this is exactly $V_i^{(d)}(t)$, \texttt{eq:localsubgraph}.
  \item \textbf{Distance-to-goal node annotation.}
    \citeauthor{knippenberg2021} includes, per visible node, ``the distance
    of each node to the goal node,'' precomputed once per instance -- this
    is exactly $\eta_i(v,t) = d_G(v,q_i(t))$, \texttt{eq:potential}.
  \item \textbf{Wait-or-move action space, shared policy.} Identical in
    both.
\end{itemize}

\subsection{Correctly, non-trivially adapted}
\begin{itemize}
  \item \textbf{Dynamic vs.\ static target.} \citeauthor{knippenberg2021}'s
    distance-to-goal values ``need only be pre-computed once for every
    problem instance,'' because each agent has exactly one job (one fixed
    goal) for the whole episode. The thesis is lifelong: an agent's target
    changes at every reassignment, and the thesis states this explicitly
    (``$\eta_i$ itself changes whenever the target changes, which happens at
    every reassignment''). This is a correct and necessary generalisation
    the source paper didn't need to make -- a good sign, not a borrowed
    shortcut.
  \item \textbf{Vertex/swap-conflict as two orthogonal equations vs.\ one
    prose definition.} \citeauthor{knippenberg2021} defines a ``node
    collision'' (two cases in prose: moving onto an occupied node, or two
    agents moving to the same node) and an ``edge collision'' (explicitly
    qualified to \emph{bi-directional} edges only). The thesis's
    \texttt{eq:vertexconflict} covers both node-collision cases with one
    inequality, and its \texttt{eq:swapconflict} needs no ``bi-directional''
    qualifier: a swap is only physically reachable when
    \texttt{eq:transition} already requires both $(v,w)\in E$ and
    $(w,v)\in E$ to hold for the two agents involved, so the bidirectionality
    falls out of constraint composition instead of being hand-added to the
    collision definition. This is a tighter formalisation than the source.
\end{itemize}

\subsection{Genuine extensions (no counterpart in the source paper)}
\begin{itemize}
  \item \textbf{The entire storage layer} -- $x_t$, $\mathcal{X}$,
    \texttt{eq:storageupdate}, \texttt{eq:storageobjective}, and the
    coupling condition \texttt{eq:coupling} that links it back to routing.
    \citeauthor{knippenberg2021} is a single-layer routing problem; this
    second, slower-timescale layer and its feedback loop is the thesis's
    actual contribution over the source paper, and is worth stating this
    plainly wherever the thesis positions itself against MAPF/MAPD
    literature.
  \item \textbf{Explicit communication graph $G_t^{\mathcal{A}}$ and
    messages.} \citeauthor{knippenberg2021} states explicitly they consider
    only ``implicit communication'' (agents observe each other's state, but
    nothing is explicitly communicated). \texttt{eq:commgraph} and the
    message component of \texttt{eq:observation} are a genuine addition.
  \item \textbf{The congestion feature $\delta_i(t)$.} No counterpart in
    the source.
  \item \textbf{Three information regimes (centralised / section-based /
    decentralised) with a shared assignment mechanism.}
    \citeauthor{knippenberg2021} studies a single decentralised regime and
    has no assignment sub-problem at all (each agent's one job is given, not
    chosen online) -- the thesis's Table~I and the assignment/routing split
    are new, driven by the lifelong setting's need to keep deciding who does
    what.
  \item \textbf{Throughput, backlog, entropy/concentration.} All specific
    to the continuous/lifelong framing; \citeauthor{knippenberg2021}'s
    objective (a single scalar per solved batch instance, eq.~1--2 in the
    source) has no analogue for an ongoing system with no end state.
\end{itemize}

\subsection{[K1, Deviation] Strong connectivity}
\citeauthor{knippenberg2021} state plainly: ``environment graphs are assumed
to be strongly connected, i.e.\ each node is reachable from all other
nodes.'' The thesis (even after [G10]'s fix) states the weaker ``$d_G(v,w)$
is finite for every pair of vertices that appears in the constructions
below.'' The source's version is simpler to state and to verify computationally
(one graph-level check, independent of what gets used later); the thesis's
version is strictly weaker and therefore more permissive of the ``irregular
blocks'' / unreachable side-room layouts its own motivating paragraph gestures
at. \textbf{This is a judgement call, not a correctness issue}: if generated
instances never actually exercise a disconnected pocket, the weaker statement
buys nothing and the source's crisper assumption should probably be adopted
outright; if partial reachability is a deliberate stress case for the
decentralised controller (an agent whose target briefly becomes unreachable),
the weaker version is doing real work and is worth keeping, explicitly framed
as a deviation from the more common assumption. The fixed .tex takes the
conservative option -- keeps the weaker statement, adds one sentence
contrasting it with \citeauthor{knippenberg2021}'s choice and naming the
tradeoff -- and leaves the actual decision to Section~VI (\emph{Instance
generation}), where it is checkable against what instances are actually
generated.

\subsection{[K2, Gap] No task deadlines}
This is the most substantial finding from the comparison. Time constraints
are \citeauthor{knippenberg2021}'s entire contribution -- their title is
literally ``Time-\emph{Constrained} MAPF'' -- realised as a job 4-tuple
$(t_i^s, t_i^d, v_i^s, v_i^g)$ with an explicit deadline $t_i^d$, and an
objective that is either a hard feasibility cut (a plan that misses its
deadline is invalid, eq.~1 in the source) or a soft penalty term
$w\cdot\max(0, t_{\pi,\text{actual}} - t_{\pi,\text{deadline}})$ added to the
cost (eq.~2). The thesis's task tuple \texttt{eq:task} carries only a release
time $r_j$; there is no deadline anywhere, and timeliness is scored only in
aggregate, after the fact, via the mean $\bar\zeta_T$ (\texttt{eq:throughput}).
Under the thesis's formulation, no individual task can be said to have
\emph{failed} -- an arbitrarily late delivery only nudges an average.

Given the thesis explicitly reuses \citeauthor{knippenberg2021}'s
architecture (built specifically to handle deadlines) while dropping the one
thing that architecture was built for, this needed to be an explicit,
stated choice rather than a silent omission -- warehouse SLAs are a real
business constraint, so a reader familiar with the source paper would
reasonably ask why deadlines didn't carry over too. \textbf{Fix (already
applied in \texttt{problem\_formulation\_fixed.tex}, Section~J):} one
paragraph stating the omission explicitly, alongside the other scope
exclusions, and naming what reintroducing deadlines would require: a
$t_j^{\mathrm{d}}$ field added to \texttt{eq:task}, plus either a feasibility
condition (hard) or a term in \texttt{eq:functional} of the same shape as
\citeauthor{knippenberg2021}'s eq.~2 (soft). Both are Problem-Formulation-level
changes, not Method-level ones, since either would redefine what a valid or
well-scored solution \emph{is} -- consistent with how the thesis already
treats \texttt{eq:coupling} and the conflict constraints as problem-level
because they define validity, not merely score it.

\subsection{[K3--K6, Deviation] Smaller, favourable deviations}
\begin{itemize}
  \item \textbf{[K3] Edge cost $c(e) > 0$ (thesis) vs.\ $c(v,v') \ge 0$
    (source).} The thesis's strict positivity rules out zero-cost cycles,
    which keeps $d_G$ well-behaved as a strictly hop-sensitive quantity and
    keeps the corrected potential-bound argument ([Error] above) free of
    degenerate edge cases. A defensible tightening, not a regression.
  \item \textbf{[K4] No \texttt{canWait}-style per-node restriction.}
    \citeauthor{knippenberg2021} give each node a boolean attribute
    controlling whether waiting is allowed there (e.g.\ a mainline track
    where a train cannot simply stop); the thesis assumes waiting is legal
    everywhere via a single $c_{\mathrm{wait}}$. This underwrites the
    thesis's ``a collision-free joint action always exists'' claim
    (Section~B) for free, but is a real loss of expressiveness relative to
    the source if a future revision wants to model genuinely
    non-stoppable segments -- worth a one-line note in Scope if that
    generality is ever wanted; not fixed here since it is not currently
    claimed as a capability.
  \item \textbf{[K5] No dummy-node start-time mechanism needed.}
    \citeauthor{knippenberg2021} materialise each agent at a private dummy
    node until its individual start time $t_i^s$ elapses, to avoid an agent
    appearing mid-graph and causing a spurious collision. The thesis avoids
    this entirely by construction: the fleet $\mathcal{A}$ is fixed and
    persistent (agents idle at $V_{ep}$ between tasks), and it is
    \emph{tasks}, not agents, that have release times $r_j$. A structural
    consequence of the lifelong framing, not a borrowed mechanism -- correctly
    absent, not simply omitted.
  \item \textbf{[K6] Reward design (Method-level, noted here for
    completeness).} \citeauthor{knippenberg2021} use a fixed table of
    empirically-tuned per-action penalties (Table~1 in the source: e.g.\
    move $-0.1$, collision $-2$, goal completion $+10$) with no stated
    invariance guarantee. \texttt{Method\_and\_Implementation.md} instead
    specifies potential-based shaping, $\gamma\Phi_i(t{+}1)-\Phi_i(t)$ with
    $\Phi_i(t) = -\eta_i(\ell_i(t),t)$, which provably preserves the
    underlying game's equilibria (\citealp{ng1999policy}; extended to
    dynamic potentials by \citealp{devlin2012dynamic}) -- directly relevant
    here since $\Phi_i$ is non-stationary by construction (the target $q_i(t)$
    changes at every reassignment). \citeauthor{knippenberg2021} even flag
    their own use of DQN as a limitation and name PPO as future work (their
    Section~8); \texttt{Method\_and\_Implementation.md}'s planned PPO
    training procedure directly answers that. Correctly scoped to Method, not
    Problem Formulation -- noted here only because the comparison is what
    surfaced it as a deliberate, citable improvement rather than an
    arbitrary choice.
\end{itemize}

\section{Where does this leave Method\_and\_Implementation.md?}
\label{sec:method-mapping}

Checking \texttt{problem\_formulation.tex} (as given, not a hypothetical
denser draft) against \texttt{Method\_and\_Implementation.md}'s own
six-item cut list:

\begin{enumerate}
  \item \emph{Controller-architecture algorithmic detail (prioritised
    planning, matching costs, zone hand-off protocols)} -- \textbf{not
    present.} Section~E as given only states the abstract information
    signature and Table~I; no algorithm-level content to cut.
  \item \emph{Encoder / message-passing equations
    ($z_i=\mathrm{Enc}_G(\cdot)$, $\phi_q,\psi_q$)} -- \textbf{not present.}
    Section~F as given defines only $V_i^{(d)}$, $\eta_i$, $\delta_i$,
    $G_t^{\mathcal{A}}$, $o_i(t)$ -- no architecture equations.
  \item \emph{Exact functional forms of $D$, $K$, $R$ and weights
    $\alpha,\beta,\chi$} -- \textbf{already deferred} in the text as given
    (``specified in Section~\ref{sec:method:storage}'').
  \item \emph{The reward function} -- \textbf{already deferred}
    (\texttt{eq:objective}'s prose: ``$R_i$ are left unspecified here'').
  \item \emph{Defensive justification passages} -- \textbf{present, low
    priority.} These exist (paraphrased in
    \texttt{Method\_and\_Implementation.md} rather than quoted verbatim) but
    are one-line footnote-style asides, not misplaced algorithmic content;
    rephrasing them is a copy-edit, not a structural fix, and wasn't applied
    here.
  \item \emph{Full notation table $\to$ appendix} -- \textbf{applies, but
    only to the table added in this review} (the original had none). Fixed:
    see [G11] above and Appendix~\ref{app:notation} below.
\end{enumerate}

The short version: \textbf{almost nothing from this review belongs in
Method\_and\_Implementation.md.} Every fix in
\texttt{problem\_formulation\_fixed.tex} is definitional (``what the object
is''), which is exactly what Section~IV is supposed to contain per that
document's own contract table. The one piece of content that needed to move
doesn't belong in Method either -- by \texttt{Method\_and\_Implementation.md}'s
own instruction, a full notation table is \emph{Appendix} material, a third
location distinct from both Problem Formulation and Method. That relocation
is done in Appendix~\ref{app:notation} of this document (ready to paste into
the thesis's own appendix chapter, which is not a file present in this
repository).

\appendix
\section{Notation disambiguation table}
\label{app:notation}

Several symbol families in Section~IV share a base letter across unrelated
objects. None are typos or the same object under different names; each pair
below denotes two genuinely different things that happen to look alike.

\begin{table}[h]
  \centering
  \caption{Symbols that share a base letter but denote unrelated objects.}
  \label{tab:notation}
  \begin{tabular}{ll}
    \toprule
    Symbols & Distinct meanings \\
    \midrule
    $c(e)$, $c_{\mathrm{wait}}$ & edge cost / wait cost (instance parameters) \\
    $\hat{c}(v,w)$ & realised one-step cost, eq.~(onestepcost) \\
    $\bar{c}_T$ & mean movement cost per completed task, eq.~(movementcost) \\
    $C_T$ & concentration index, eq.~(concentration) \\
    $\mathcal{C}_t$ & completed-task set, eq.~(lifecycle) \\
    \addlinespace
    $\mathcal{Q}_t$ & waiting-task set, eq.~(lifecycle) \\
    $q_i(t)$ & agent $i$'s current target vertex, eq.~(potential) \\
    $q_{\mathrm{svc}}, q_{\mathrm{mov}}, q_{\mathrm{wait}}, q_{\mathrm{conc}}, q_{\mathrm{back}}$
      & weights in the system-level criterion, eq.~(functional) \\
    \addlinespace
    $\mathcal{O}_i$ & observation space, eq.~(posg) \\
    $O_i$ & observation function, eq.~(posg) \\
    \addlinespace
    $\mathcal{U}(v)$ & action set at vertex $v$, eq.~(actions) \\
    $\mathcal{U}_i = \bigcup_{v} \mathcal{U}(v)$ & agent $i$'s fixed action space in $\mathcal{M}$, eq.~(posg) \\
    \bottomrule
  \end{tabular}
\end{table}

\noindent When pasted into the thesis's actual appendix, replace each
\texttt{eq.~(xxx)} above with \verb|\eqref{eq:xxx}| -- the labels already
match \texttt{problem\_formulation\_fixed.tex} verbatim, so this table can be
dropped in without renumbering anything.

\begin{thebibliography}{9}
\bibitem[van Knippenberg et~al.(2021)]{knippenberg2021}
Marijn van Knippenberg, Mike Holenderski, and Vlado Menkovski.
\newblock Time-Constrained Multi-Agent Path Finding in Non-Lattice Graphs
  with Deep Reinforcement Learning.
\newblock In \emph{Proceedings of the 13th Asian Conference on Machine
  Learning}, volume 157 of \emph{PMLR}, 2021.

\bibitem[Ng et~al.(1999)]{ng1999policy}
Andrew~Y. Ng, Daishi Harada, and Stuart Russell.
\newblock Policy Invariance Under Reward Transformations: Theory and
  Application to Reward Shaping.
\newblock In \emph{ICML}, 1999.

\bibitem[Devlin and Kudenko(2012)]{devlin2012dynamic}
Sam~Michael Devlin and Daniel Kudenko.
\newblock Dynamic Potential-Based Reward Shaping.
\newblock In \emph{AAMAS}, 2012.
\end{thebibliography}

\end{document}
